\section{Reconstruct green and the two residuals}
LUMA is a two-byte command. Its first byte has six bits for the green change with a
bias of thirty-two. The high and low nibbles of the second byte contain red-minus-green
and blue-minus-green changes, each biased by eight. The name does not mean that the
decoder computes an approximate luminance or discards chroma. The representation is
lossless because both residuals are retained.

Let \(d_g\) be the first field minus thirty-two, \(u\) the high nibble minus eight,
and \(v\) the low nibble minus eight. Then
\[
  d_r=d_g+u,\qquad d_b=d_g+v.
\]
Apply these changes to the previous red, green and blue values modulo 256, retaining
alpha. The green change ranges from minus thirty-two through thirty-one. Each residual
ranges from minus eight through seven. Those ranges differ; do not use the same mask
or bias for all three fields.

The bytes \texttt{a5 97} give green change five, red residual one and blue residual
minus one. Their complete changes are therefore six, five and four. From
\((1,1,1,255)\), reconstruct \((7,6,5,255)\). Reading the second byte as a signed
integer would obscure the two independent nibbles. Load it unsigned, then split it.

An encoder tests these ranges only when alpha is unchanged. If any required change
does not fit, RGB can still represent the three new components. If alpha changes,
use RGBA unless an already cached complete pixel provides an INDEX representation.
